% Options for packages loaded elsewhere
\PassOptionsToPackage{unicode}{hyperref}
\PassOptionsToPackage{hyphens}{url}
\documentclass[
]{article}
\usepackage{xcolor}
\usepackage{amsmath,amssymb}
\setcounter{secnumdepth}{5}
\usepackage{iftex}
\ifPDFTeX
  \usepackage[T1]{fontenc}
  \usepackage[utf8]{inputenc}
  \usepackage{textcomp} % provide euro and other symbols
\else % if luatex or xetex
  \usepackage{unicode-math} % this also loads fontspec
  \defaultfontfeatures{Scale=MatchLowercase}
  \defaultfontfeatures[\rmfamily]{Ligatures=TeX,Scale=1}
\fi
\usepackage{lmodern}
\ifPDFTeX\else
  % xetex/luatex font selection
\fi
% Use upquote if available, for straight quotes in verbatim environments
\IfFileExists{upquote.sty}{\usepackage{upquote}}{}
\IfFileExists{microtype.sty}{% use microtype if available
  \usepackage[]{microtype}
  \UseMicrotypeSet[protrusion]{basicmath} % disable protrusion for tt fonts
}{}
\makeatletter
\@ifundefined{KOMAClassName}{% if non-KOMA class
  \IfFileExists{parskip.sty}{%
    \usepackage{parskip}
  }{% else
    \setlength{\parindent}{0pt}
    \setlength{\parskip}{6pt plus 2pt minus 1pt}}
}{% if KOMA class
  \KOMAoptions{parskip=half}}
\makeatother
\usepackage{longtable,booktabs,array}
\newcounter{none} % for unnumbered tables
\usepackage{calc} % for calculating minipage widths
% Correct order of tables after \paragraph or \subparagraph
\usepackage{etoolbox}
\makeatletter
\patchcmd\longtable{\par}{\if@noskipsec\mbox{}\fi\par}{}{}
\makeatother
% Allow footnotes in longtable head/foot
\IfFileExists{footnotehyper.sty}{\usepackage{footnotehyper}}{\usepackage{footnote}}
\makesavenoteenv{longtable}
\setlength{\emergencystretch}{3em} % prevent overfull lines
\providecommand{\tightlist}{%
  \setlength{\itemsep}{0pt}\setlength{\parskip}{0pt}}
\usepackage{bookmark}
\IfFileExists{xurl.sty}{\usepackage{xurl}}{} % add URL line breaks if available
\urlstyle{same}
\hypersetup{
  hidelinks,
  pdfcreator={LaTeX via pandoc}}

\author{}
\date{}

\begin{document}

{
\setcounter{tocdepth}{3}
\tableofcontents
}
Yes. I read the full B2 Advanced Math material. I'm going to teach this
as a \textbf{complete Digital SAT Advanced Math lesson}, not as notes to
memorize. The domain is built around four official skills: equivalent
expressions, nonlinear equations in one variable, nonlinear systems, and
nonlinear functions.

The six question families your file identifies as especially common are
equivalent-expression problems, quadratic vertex/min-max problems,
exponential growth/decay, extraneous-solution questions, nonlinear
systems, and rearranging nonlinear formulas.

\begin{center}\rule{0.5\linewidth}{0.5pt}\end{center}

\section{0. The mental model for the whole SAT Advanced Math
domain}\label{the-mental-model-for-the-whole-sat-advanced-math-domain}

Almost every problem in this domain is asking you to do one of these:

{\def\LTcaptype{none} % do not increment counter
\begin{longtable}[]{@{}
  >{\raggedright\arraybackslash}p{(\linewidth - 2\tabcolsep) * \real{0.5000}}
  >{\raggedright\arraybackslash}p{(\linewidth - 2\tabcolsep) * \real{0.5000}}@{}}
\toprule\noalign{}
\begin{minipage}[b]{\linewidth}\raggedright
SAT task
\end{minipage} & \begin{minipage}[b]{\linewidth}\raggedright
What you are actually doing
\end{minipage} \\
\midrule\noalign{}
\endhead
\bottomrule\noalign{}
\endlastfoot
``Which expression is equivalent?'' & Rewrite without changing value \\
``Solve for \(x\)'' & Find inputs that make an equation true \\
``How many solutions?'' & Count intersections/roots without necessarily
solving everything \\
``What is the minimum/maximum?'' & Find a vertex \\
``What does this constant represent?'' & Interpret parameters in a
model \\
``Which graph could represent\ldots?'' & Connect algebraic form to graph
features \\
``For what value of \(k\)\ldots?'' & Use discriminants, tangency, or
other constraints \\
``Solve for \(r\), \(t\), etc.'' & Rearrange a formula \\
``Which function models\ldots?'' & Translate words into
quadratic/exponential/etc. form \\
\end{longtable}
}

Before doing algebra, \textbf{name the type of problem}.

That alone prevents a lot of wasted work.

\begin{center}\rule{0.5\linewidth}{0.5pt}\end{center}

\section{PART I. FOUNDATIONAL FORMULAS YOU MUST
KNOW}\label{part-i.-foundational-formulas-you-must-know}

\section{1. Exponent rules}\label{exponent-rules}

These need to become automatic.

{\def\LTcaptype{none} % do not increment counter
\begin{longtable}[]{@{}
  >{\raggedright\arraybackslash}p{(\linewidth - 4\tabcolsep) * \real{0.3333}}
  >{\raggedright\arraybackslash}p{(\linewidth - 4\tabcolsep) * \real{0.3333}}
  >{\raggedright\arraybackslash}p{(\linewidth - 4\tabcolsep) * \real{0.3333}}@{}}
\toprule\noalign{}
\begin{minipage}[b]{\linewidth}\raggedright
Situation
\end{minipage} & \begin{minipage}[b]{\linewidth}\raggedright
Rule
\end{minipage} & \begin{minipage}[b]{\linewidth}\raggedright
Example
\end{minipage} \\
\midrule\noalign{}
\endhead
\bottomrule\noalign{}
\endlastfoot
Same base, multiply & \(a^m a^n=a^{m+n}\) & \(x^3x^5=x^8\) \\
Same base, divide & \(\frac{a^m}{a^n}=a^{m-n}\) &
\(\frac{x^7}{x^2}=x^5\) \\
Power of power & \((a^m)^n=a^{mn}\) & \((x^3)^4=x^{12}\) \\
Power of product & \((ab)^n=a^nb^n\) & \((2x)^3=8x^3\) \\
Zero exponent & \(a^0=1\) & \(7^0=1\) \\
Negative exponent & \(a^{-n}=\frac1{a^n}\) & \(x^{-3}=\frac1{x^3}\) \\
Fraction exponent & \(a^{m/n}=\sqrt[n]{a^m}\) &
\(x^{2/3}=\sqrt[3]{x^2}\) \\
\end{longtable}
}

\subsection{Important distinction}\label{important-distinction}

Multiplication lets you \textbf{add exponents} only when the bases
match:

\[
x^{1/2}x^{1/3}
=
x^{1/2+1/3}
=
x^{5/6}.
\]

Not \(x^{1/6}\).

For a power raised to another power, multiply:

\[
(x^{2/3})^6=x^4.
\]

\subsection{SAT example}\label{sat-example}

Simplify

\[
\frac{x^{7/3}}{x^{1/2}}.
\]

Subtract:

\[
x^{7/3-1/2}
=
x^{14/6-3/6}
=
x^{11/6}.
\]

\begin{center}\rule{0.5\linewidth}{0.5pt}\end{center}

\section{2. Radical rules}\label{radical-rules}

A radical and a fractional exponent are two ways of writing the same
thing:

\[
\sqrt{x}=x^{1/2}
\]

\[
\sqrt[3]{x}=x^{1/3}
\]

\[
\sqrt[4]{x^3}=x^{3/4}.
\]

This is extremely useful because complicated radical expressions become
ordinary exponent problems.

\subsection{Critical exception}\label{critical-exception}

\[
\boxed{\sqrt{x^2}=|x|}
\]

not automatically \(x\).

Why?

If \(x=-5\),

\[
\sqrt{x^2}=\sqrt{25}=5,
\]

not \(-5\).

\begin{center}\rule{0.5\linewidth}{0.5pt}\end{center}

\section{3. Rationalizing
denominators}\label{rationalizing-denominators}

For something simple:

\[
\frac1{\sqrt 5}
\]

multiply numerator and denominator by \(\sqrt5\):

\[
\frac{\sqrt5}{5}.
\]

For a binomial denominator:

\[
\frac1{3+\sqrt2},
\]

use the \textbf{conjugate}:

\[
3-\sqrt2.
\]

Then

\[
\frac{1}{3+\sqrt2}
\cdot
\frac{3-\sqrt2}{3-\sqrt2}
=
\frac{3-\sqrt2}{9-2}
=
\frac{3-\sqrt2}{7}.
\]

This works because

\[
(a+b)(a-b)=a^2-b^2.
\]

\begin{center}\rule{0.5\linewidth}{0.5pt}\end{center}

\section{PART II. EQUIVALENT
EXPRESSIONS}\label{part-ii.-equivalent-expressions}

Your file makes this the first major Advanced Math section. Two
expressions are equivalent when they produce the same value for every
allowable input.

The core operations are:

expanding, factoring, simplifying rational expressions, manipulating
exponents/radicals, completing the square, and switching quadratic
forms.

\begin{center}\rule{0.5\linewidth}{0.5pt}\end{center}

\section{4. Polynomial addition and
subtraction}\label{polynomial-addition-and-subtraction}

You may combine only \textbf{like terms}.

\[
3x^2+5x^2=8x^2
\]

but

\[
3x^2+5x
\]

cannot combine.

\subsection{Subtraction trap}\label{subtraction-trap}

Suppose:

\[
(4x^2+3x-2)-(x^2-5x+7).
\]

Distribute the negative to \textbf{every term}:

\[
4x^2+3x-2-x^2+5x-7.
\]

Then:

\[
3x^2+8x-9.
\]

The SAT loves the sign error

\[
-(x^2-5x+7)
\]

being incorrectly written as

\[
-x^2-5x-7.
\]

Correct:

\[
-x^2+5x-7.
\]

\begin{center}\rule{0.5\linewidth}{0.5pt}\end{center}

\section{5. Multiplying polynomials}\label{multiplying-polynomials}

Every term must multiply every relevant term.

Example:

\[
(2x-3)(x^2+4x-5).
\]

Multiply \(2x\):

\[
2x^3+8x^2-10x.
\]

Multiply \(-3\):

\[
-3x^2-12x+15.
\]

Combine:

\[
\boxed{2x^3+5x^2-22x+15}.
\]

\begin{center}\rule{0.5\linewidth}{0.5pt}\end{center}

\section{6. Three identities you should recognize
instantly}\label{three-identities-you-should-recognize-instantly}

\subsection{Square of a sum}\label{square-of-a-sum}

\[
\boxed{(a+b)^2=a^2+2ab+b^2}
\]

Example:

\[
(x+7)^2=x^2+14x+49.
\]

\subsection{Square of a difference}\label{square-of-a-difference}

\[
\boxed{(a-b)^2=a^2-2ab+b^2}
\]

Example:

\[
(3x-2)^2
=
9x^2-12x+4.
\]

\subsection{Difference of squares}\label{difference-of-squares}

\[
\boxed{a^2-b^2=(a-b)(a+b)}
\]

Example:

\[
x^2-81=(x-9)(x+9).
\]

\subsubsection{Huge trap}\label{huge-trap}

\[
(a+b)^2 \neq a^2+b^2.
\]

The middle term \(2ab\) must exist.

\begin{center}\rule{0.5\linewidth}{0.5pt}\end{center}

\section{7. Factoring: the master
method}\label{factoring-the-master-method}

Factoring means reversing multiplication.

Your decision process should be:

\textbf{GCF first → count terms → identify pattern → factor again if
possible.}

\subsection{7.1 Greatest common factor}\label{greatest-common-factor}

\[
12x^3-8x^2+4x.
\]

Everything contains \(4x\):

\[
\boxed{4x(3x^2-2x+1)}.
\]

Always check GCF first.

\begin{center}\rule{0.5\linewidth}{0.5pt}\end{center}

\section{7.2 Difference of squares}\label{difference-of-squares-1}

\[
x^2-25
\]

is

\[
x^2-5^2
\]

so

\[
\boxed{(x-5)(x+5)}.
\]

More difficult:

\[
x^4-16.
\]

First:

\[
(x^2-4)(x^2+4).
\]

But \(x^2-4\) factors again:

\[
\boxed{(x-2)(x+2)(x^2+4)}.
\]

Do not stop too early.

\begin{center}\rule{0.5\linewidth}{0.5pt}\end{center}

\section{7.3 Sum of squares does NOT factor over the
reals}\label{sum-of-squares-does-not-factor-over-the-reals}

\[
x^2+25
\]

does not factor into real linear factors.

This distinction matters:

\[
x^2-25 \quad \checkmark
\]

\[
x^2+25 \quad \text{no real factorization}.
\]

\begin{center}\rule{0.5\linewidth}{0.5pt}\end{center}

\section{7.4 Perfect-square trinomials}\label{perfect-square-trinomials}

Recognize:

\[
a^2+2ab+b^2=(a+b)^2
\]

and

\[
a^2-2ab+b^2=(a-b)^2.
\]

Example:

\[
x^2+10x+25.
\]

First and last:

\[
x^2,\quad 25=5^2.
\]

Middle:

\[
2(x)(5)=10x.
\]

Therefore:

\[
\boxed{(x+5)^2}.
\]

Example:

\[
4x^2-12x+9.
\]

\[
(2x)^2,\quad 3^2,
\]

middle:

\[
-2(2x)(3)=-12x.
\]

So

\[
\boxed{(2x-3)^2}.
\]

\begin{center}\rule{0.5\linewidth}{0.5pt}\end{center}

\section{7.5 Simple trinomials}\label{simple-trinomials}

For

\[
x^2+bx+c,
\]

find two numbers \(p,q\) such that

\[
p+q=b
\]

and

\[
pq=c.
\]

Example:

\[
x^2-7x+12.
\]

Need:

\[
p+q=-7,\qquad pq=12.
\]

Numbers:

\[
-3,-4.
\]

Therefore:

\[
\boxed{(x-3)(x-4)}.
\]

\begin{center}\rule{0.5\linewidth}{0.5pt}\end{center}

\section{\texorpdfstring{7.6 Hard trinomials:
\(ax^2+bx+c\)}{7.6 Hard trinomials: ax\^{}2+bx+c}}\label{hard-trinomials-ax2bxc}

Example:

\[
6x^2+13x-5.
\]

Multiply:

\[
ac=6(-5)=-30.
\]

Find two numbers multiplying to \(-30\) and adding to \(13\):

\[
15,-2.
\]

Split the middle:

\[
6x^2+15x-2x-5.
\]

Group:

\[
3x(2x+5)-1(2x+5).
\]

Factor common binomial:

\[
\boxed{(3x-1)(2x+5)}.
\]

This method is very reliable.

\begin{center}\rule{0.5\linewidth}{0.5pt}\end{center}

\section{7.7 Grouping}\label{grouping}

Usually four terms.

Example:

\[
x^3+3x^2-4x-12.
\]

Group:

\[
(x^3+3x^2)+(-4x-12).
\]

Factor:

\[
x^2(x+3)-4(x+3).
\]

Common factor:

\[
(x+3)(x^2-4).
\]

Difference of squares:

\[
\boxed{(x+3)(x-2)(x+2)}.
\]

\begin{center}\rule{0.5\linewidth}{0.5pt}\end{center}

\section{7.8 Cubes}\label{cubes}

Know:

\[
\boxed{a^3-b^3=(a-b)(a^2+ab+b^2)}
\]

\[
\boxed{a^3+b^3=(a+b)(a^2-ab+b^2)}.
\]

Mnemonic:

\textbf{same, opposite, always positive.}

For \(a^3-b^3\):

first factor uses same sign:

\[
a-b.
\]

Middle quadratic uses opposite sign:

\[
+a b.
\]

Final term positive.

Example:

\[
27x^3-64
=
(3x)^3-4^3
\]

so

\[
\boxed{(3x-4)(9x^2+12x+16)}.
\]

\begin{center}\rule{0.5\linewidth}{0.5pt}\end{center}

\section{8. Rational expressions}\label{rational-expressions}

Example:

\[
\frac{x^2-9}{x^2+5x+6}.
\]

Factor:

\[
x^2-9=(x-3)(x+3)
\]

\[
x^2+5x+6=(x+2)(x+3).
\]

Thus:

\[
\frac{(x-3)(x+3)}
{(x+2)(x+3)}
=
\boxed{\frac{x-3}{x+2}}.
\]

But the original denominator told us:

\[
x\neq -3,-2.
\]

Even though \(x+3\) disappeared algebraically, \(x=-3\) remains
excluded.

\begin{center}\rule{0.5\linewidth}{0.5pt}\end{center}

\section{9. You cancel factors, not
terms}\label{you-cancel-factors-not-terms}

Legal:

\[
\frac{(x+2)(x-5)}{(x+2)(x+3)}
\]

because \(x+2\) is a complete factor.

Illegal:

\[
\frac{x+3}{x}
\]

and trying to cancel the \(x\).

Why?

The numerator is a \textbf{sum}.

There is no common multiplicative factor of \(x\).

\begin{center}\rule{0.5\linewidth}{0.5pt}\end{center}

\section{10. Adding rational
expressions}\label{adding-rational-expressions}

You need a common denominator.

Example:

\[
\frac{2}{x}+\frac{3}{x+1}.
\]

LCD:

\[
x(x+1).
\]

Rewrite:

\[
\frac{2(x+1)}{x(x+1)}
+
\frac{3x}{x(x+1)}.
\]

Combine:

\[
\boxed{\frac{5x+2}{x(x+1)}}.
\]

Restrictions:

\[
x\neq0,-1.
\]

\begin{center}\rule{0.5\linewidth}{0.5pt}\end{center}

\section{11. Dividing rational
expressions}\label{dividing-rational-expressions}

Division means multiply by reciprocal.

\[
\frac{x^2-4}{x+1}
\div
\frac{x-2}{x+3}.
\]

Rewrite:

\[
\frac{(x-2)(x+2)}{x+1}
\cdot
\frac{x+3}{x-2}.
\]

Cancel:

\[
\boxed{\frac{(x+2)(x+3)}{x+1}}.
\]

\begin{center}\rule{0.5\linewidth}{0.5pt}\end{center}

\section{12. Radical and exponent
rewriting}\label{radical-and-exponent-rewriting}

Example:

\[
\frac{\sqrt[3]{x^2}}{\sqrt{x}}.
\]

Convert:

\[
\frac{x^{2/3}}{x^{1/2}}
=
x^{2/3-1/2}.
\]

Common denominator \(6\):

\[
\frac46-\frac36=\frac16.
\]

Therefore:

\[
\boxed{x^{1/6}}
=
\boxed{\sqrt[6]x}.
\]

This fractional-exponent strategy is usually the cleanest SAT method.

\begin{center}\rule{0.5\linewidth}{0.5pt}\end{center}

\section{PART III. COMPLETING THE
SQUARE}\label{part-iii.-completing-the-square}

\section{13. Why completing the square
exists}\label{why-completing-the-square-exists}

You are converting

\[
ax^2+bx+c
\]

into

\[
a(x-h)^2+k.
\]

The second form exposes the \textbf{vertex} immediately.

For

\[
x^2+bx,
\]

the magic number is

\[
\left(\frac b2\right)^2.
\]

\begin{center}\rule{0.5\linewidth}{0.5pt}\end{center}

\section{\texorpdfstring{14. Example with
\(a=1\)}{14. Example with a=1}}\label{example-with-a1}

Rewrite:

\[
x^2+8x+3.
\]

Half of \(8\):

\[
4.
\]

Square it:

\[
16.
\]

Add and subtract:

\[
x^2+8x+16-16+3.
\]

Perfect square:

\[
(x+4)^2-13.
\]

So:

\[
\boxed{x^2+8x+3=(x+4)^2-13}.
\]

Vertex:

\[
\boxed{(-4,-13)}.
\]

\begin{center}\rule{0.5\linewidth}{0.5pt}\end{center}

\section{\texorpdfstring{15. Example where
\(a\neq1\)}{15. Example where a\textbackslash neq1}}\label{example-where-aneq1}

Rewrite:

\[
2x^2-12x+5.
\]

Factor \(2\) from only the \(x\)-terms:

\[
2(x^2-6x)+5.
\]

Half of \(-6\):

\[
-3.
\]

Square:

\[
9.
\]

\[
2(x^2-6x+9-9)+5.
\]

\[
2[(x-3)^2-9]+5.
\]

Distribute the \(2\):

\[
2(x-3)^2-18+5.
\]

\[
\boxed{2(x-3)^2-13}.
\]

Vertex:

\[
(3,-13).
\]

\subsubsection{Main trap}\label{main-trap}

You cannot forget that the \(-9\) inside gets multiplied by \(2\).

\begin{center}\rule{0.5\linewidth}{0.5pt}\end{center}

\section{PART IV. THE THREE FORMS OF A
QUADRATIC}\label{part-iv.-the-three-forms-of-a-quadratic}

This is one of the highest-yield SAT concepts.

\section{16. Standard form}\label{standard-form}

\[
\boxed{f(x)=ax^2+bx+c}
\]

Immediately gives:

\[
y\text{-intercept}=c.
\]

Because:

\[
f(0)=c.
\]

Example:

\[
f(x)=3x^2-7x+12.
\]

\(y\)-intercept:

\[
\boxed{12}.
\]

\begin{center}\rule{0.5\linewidth}{0.5pt}\end{center}

\section{17. Factored form}\label{factored-form}

\[
\boxed{f(x)=a(x-r_1)(x-r_2)}
\]

Immediately gives roots:

\[
x=r_1,\qquad x=r_2.
\]

Example:

\[
f(x)=2(x-3)(x+5).
\]

Roots:

\[
\boxed{3,-5}.
\]

Axis of symmetry:

\[
x=\frac{r_1+r_2}{2}
\]

so:

\[
x=\frac{3+(-5)}2=-1.
\]

\begin{center}\rule{0.5\linewidth}{0.5pt}\end{center}

\section{18. Vertex form}\label{vertex-form}

\[
\boxed{f(x)=a(x-h)^2+k}
\]

Immediately gives:

vertex:

\[
\boxed{(h,k)}
\]

axis:

\[
\boxed{x=h}
\]

and min/max value:

\[
\boxed{k}.
\]

If \(a>0\), opens up → minimum.

If \(a<0\), opens down → maximum.

Example:

\[
f(x)=-4(x-2)^2+9.
\]

Vertex:

\[
(2,9).
\]

Since \(a=-4<0\):

\[
\boxed{\text{maximum}=9}.
\]

\begin{center}\rule{0.5\linewidth}{0.5pt}\end{center}

\section{19. Know which form to use}\label{know-which-form-to-use}

{\def\LTcaptype{none} % do not increment counter
\begin{longtable}[]{@{}ll@{}}
\toprule\noalign{}
Question asks for & Best form \\
\midrule\noalign{}
\endhead
\bottomrule\noalign{}
\endlastfoot
\(y\)-intercept & Standard \\
roots/zeros/\(x\)-intercepts & Factored \\
vertex & Vertex \\
maximum/minimum & Vertex \\
axis of symmetry & Vertex or factored \\
equivalent expression & Usually expand/factor \\
graph behavior & Depends on feature \\
\end{longtable}
}

A major SAT skill is not doing unnecessary algebra.

If you see:

\[
g(x)=3(x+7)^2-11
\]

and the question asks for the minimum, you should immediately answer:

\[
\boxed{-11}.
\]

\begin{center}\rule{0.5\linewidth}{0.5pt}\end{center}

\section{PART V. SOLVING QUADRATIC
EQUATIONS}\label{part-v.-solving-quadratic-equations}

\section{20. First rule: put it equal to
zero}\label{first-rule-put-it-equal-to-zero}

Quadratic equation:

\[
ax^2+bx+c=0.
\]

Then choose a method.

\begin{center}\rule{0.5\linewidth}{0.5pt}\end{center}

\section{21. Quadratic method decision
tree}\label{quadratic-method-decision-tree}

{\def\LTcaptype{none} % do not increment counter
\begin{longtable}[]{@{}ll@{}}
\toprule\noalign{}
Situation & Best method \\
\midrule\noalign{}
\endhead
\bottomrule\noalign{}
\endlastfoot
\(x^2=k\) or \((x-h)^2=k\) & Square root \\
Factors obvious & Factoring \\
Need vertex & Complete square \\
Messy / uncertain & Quadratic formula \\
Need only number of solutions & Discriminant \\
Need decimal/intersection & Desmos can help \\
\end{longtable}
}

\begin{center}\rule{0.5\linewidth}{0.5pt}\end{center}

\section{22. Factoring to solve}\label{factoring-to-solve}

\[
x^2-5x+6=0.
\]

Factor:

\[
(x-2)(x-3)=0.
\]

Zero-product property:

\[
x-2=0
\]

or

\[
x-3=0.
\]

Therefore:

\[
\boxed{x=2,3}.
\]

\begin{center}\rule{0.5\linewidth}{0.5pt}\end{center}

\section{23. Square-root method}\label{square-root-method}

\[
3(x-4)^2=27.
\]

Divide by \(3\):

\[
(x-4)^2=9.
\]

Take square root:

\[
x-4=\pm3.
\]

So:

\[
\boxed{x=7,1}.
\]

\subsubsection{Essential trap}\label{essential-trap}

\[
x^2=25
\]

means

\[
\boxed{x=\pm5},
\]

not just \(5\).

\begin{center}\rule{0.5\linewidth}{0.5pt}\end{center}

\section{24. Quadratic formula}\label{quadratic-formula}

Memorize:

\[
\boxed{x=\frac{-b\pm\sqrt{b^2-4ac}}{2a}}.
\]

Example:

\[
2x^2-7x=3.
\]

Move everything:

\[
2x^2-7x-3=0.
\]

So:

\[
a=2,\quad b=-7,\quad c=-3.
\]

Substitute:

\[
x=
\frac{-(-7)\pm\sqrt{(-7)^2-4(2)(-3)}}{4}.
\]

\[
=
\frac{7\pm\sqrt{49+24}}4.
\]

\[
\boxed{x=\frac{7\pm\sqrt{73}}4}.
\]

Write down \(a,b,c\) first. It dramatically reduces sign errors.

\begin{center}\rule{0.5\linewidth}{0.5pt}\end{center}

\section{25. The discriminant}\label{the-discriminant}

Inside the radical:

\[
\boxed{D=b^2-4ac}.
\]

It tells you the number of \textbf{real} roots.

{\def\LTcaptype{none} % do not increment counter
\begin{longtable}[]{@{}lr@{}}
\toprule\noalign{}
Discriminant & Real roots \\
\midrule\noalign{}
\endhead
\bottomrule\noalign{}
\endlastfoot
\(D>0\) & 2 \\
\(D=0\) & 1 \\
\(D<0\) & 0 \\
\end{longtable}
}

Geometrically:

{\def\LTcaptype{none} % do not increment counter
\begin{longtable}[]{@{}ll@{}}
\toprule\noalign{}
\(D\) & Parabola \\
\midrule\noalign{}
\endhead
\bottomrule\noalign{}
\endlastfoot
\(>0\) & crosses \(x\)-axis twice \\
\(=0\) & touches \(x\)-axis \\
\(<0\) & never reaches \(x\)-axis \\
\end{longtable}
}

\begin{center}\rule{0.5\linewidth}{0.5pt}\end{center}

\section{26. Parameter questions using
discriminant}\label{parameter-questions-using-discriminant}

Example:

For what \(k\) does

\[
x^2+kx+9=0
\]

have exactly one real solution?

Exactly one means:

\[
D=0.
\]

\[
k^2-4(1)(9)=0.
\]

\[
k^2-36=0.
\]

\[
k^2=36.
\]

\[
\boxed{k=\pm6}.
\]

Do not forget both signs.

\begin{center}\rule{0.5\linewidth}{0.5pt}\end{center}

\section{27. Two distinct solutions}\label{two-distinct-solutions}

Suppose:

\[
2x^2-4x+k=0
\]

has two distinct real roots.

Need:

\[
D>0.
\]

\[
(-4)^2-4(2)(k)>0.
\]

\[
16-8k>0.
\]

\[
16>8k.
\]

\[
\boxed{k<2}.
\]

\begin{center}\rule{0.5\linewidth}{0.5pt}\end{center}

\section{PART VI. QUADRATIC-IN-FORM
PROBLEMS}\label{part-vi.-quadratic-in-form-problems}

\section{\texorpdfstring{28.
\(u\)-substitution}{28. u-substitution}}\label{u-substitution}

Sometimes an equation is not technically quadratic in \(x\), but
contains a repeated expression.

Example:

\[
x^4-5x^2+6=0.
\]

Notice:

\[
x^4=(x^2)^2.
\]

Let:

\[
u=x^2.
\]

Then:

\[
u^2-5u+6=0.
\]

Factor:

\[
(u-2)(u-3)=0.
\]

So:

\[
u=2,3.
\]

Back-substitute:

\[
x^2=2
\]

or

\[
x^2=3.
\]

Thus:

\[
\boxed{x=\pm\sqrt2,\pm\sqrt3}.
\]

Four real solutions.

\begin{center}\rule{0.5\linewidth}{0.5pt}\end{center}

\section{29. Negative substituted
values}\label{negative-substituted-values}

Example:

\[
x^4-2x^2-3=0.
\]

Set \(u=x^2\):

\[
u^2-2u-3=0.
\]

Factor:

\[
(u-3)(u+1)=0.
\]

So:

\[
u=3,-1.
\]

Back-substitute:

\[
x^2=3 \Rightarrow x=\pm\sqrt3.
\]

But:

\[
x^2=-1
\]

has no real solution.

Therefore:

\[
\boxed{2\text{ real solutions}}.
\]

\begin{center}\rule{0.5\linewidth}{0.5pt}\end{center}

\section{PART VII. RADICAL EQUATIONS}\label{part-vii.-radical-equations}

\section{30. Why extraneous solutions
appear}\label{why-extraneous-solutions-appear}

Squaring is not a one-to-one operation.

For example:

\[
3^2=(-3)^2.
\]

Therefore when you square both sides of an equation, information about
sign can disappear.

So you must check answers in the \textbf{original equation}.

\begin{center}\rule{0.5\linewidth}{0.5pt}\end{center}

\section{31. Radical equation method}\label{radical-equation-method}

Suppose:

\[
\sqrt{2x+3}+5=x.
\]

Isolate the radical:

\[
\sqrt{2x+3}=x-5.
\]

Square:

\[
2x+3=(x-5)^2.
\]

\[
2x+3=x^2-10x+25.
\]

Move everything:

\[
x^2-12x+22=0.
\]

Quadratic formula gives:

\[
x=6\pm\sqrt{14}.
\]

Now check.

Since the original isolated form was

\[
\sqrt{2x+3}=x-5,
\]

the right side must be nonnegative.

For

\[
x=6-\sqrt{14}\approx2.26,
\]

\[
x-5<0.
\]

Impossible because a square root is nonnegative.

Reject it.

Therefore:

\[
\boxed{x=6+\sqrt{14}}.
\]

\begin{center}\rule{0.5\linewidth}{0.5pt}\end{center}

\section{32. Radical SAT checklist}\label{radical-sat-checklist}

For any square-root equation:

\[
\boxed{\text{isolate → square → solve → check original}}
\]

Never skip the final step.

\begin{center}\rule{0.5\linewidth}{0.5pt}\end{center}

\section{PART VIII. RATIONAL
EQUATIONS}\label{part-viii.-rational-equations}

\section{33. Start with restrictions}\label{start-with-restrictions}

Example:

\[
\frac{x^2}{x-1}=\frac1{x-1}.
\]

Original denominator means:

\[
\boxed{x\neq1}.
\]

Now multiply by \(x-1\):

\[
x^2=1.
\]

So:

\[
x=\pm1.
\]

But \(x=1\) is illegal.

Thus:

\[
\boxed{x=-1}.
\]

\begin{center}\rule{0.5\linewidth}{0.5pt}\end{center}

\section{34. Full rational-equation
process}\label{full-rational-equation-process}

For:

\[
\frac{x}{x-3}
+
\frac2{x+1}
=
\frac5{x^2-2x-3},
\]

factor:

\[
x^2-2x-3=(x-3)(x+1).
\]

Restrictions:

\[
x\neq3,-1.
\]

LCD:

\[
(x-3)(x+1).
\]

Multiply everything:

\[
x(x+1)+2(x-3)=5.
\]

Expand:

\[
x^2+x+2x-6=5.
\]

\[
x^2+3x-11=0.
\]

Then solve the quadratic.

Final candidates must still respect:

\[
x\neq3,-1.
\]

\begin{center}\rule{0.5\linewidth}{0.5pt}\end{center}

\section{PART IX. ABSOLUTE VALUE
EQUATIONS}\label{part-ix.-absolute-value-equations}

\section{35. Meaning of absolute value}\label{meaning-of-absolute-value}

\[
|x|
\]

is the distance between \(x\) and \(0\).

Distances cannot be negative.

Therefore:

\[
|x|=-5
\]

has

\[
\boxed{\text{no solution}}.
\]

\begin{center}\rule{0.5\linewidth}{0.5pt}\end{center}

\section{36. Standard absolute-value
equation}\label{standard-absolute-value-equation}

If:

\[
|A|=B
\]

and \(B>0\), then:

\[
A=B
\]

or

\[
A=-B.
\]

Example:

\[
|2x+5|-3=10.
\]

Isolate:

\[
|2x+5|=13.
\]

Case 1:

\[
2x+5=13
\]

\[
2x=8
\]

\[
x=4.
\]

Case 2:

\[
2x+5=-13
\]

\[
2x=-18
\]

\[
x=-9.
\]

Answer:

\[
\boxed{4,-9}.
\]

\begin{center}\rule{0.5\linewidth}{0.5pt}\end{center}

\section{37. Number of absolute-value
solutions}\label{number-of-absolute-value-solutions}

For:

\[
|x-h|=k,
\]

{\def\LTcaptype{none} % do not increment counter
\begin{longtable}[]{@{}lr@{}}
\toprule\noalign{}
\(k\) & Number of solutions \\
\midrule\noalign{}
\endhead
\bottomrule\noalign{}
\endlastfoot
\(k>0\) & 2 \\
\(k=0\) & 1 \\
\(k<0\) & 0 \\
\end{longtable}
}

Visualize a horizontal line intersecting a V.

\begin{center}\rule{0.5\linewidth}{0.5pt}\end{center}

\section{PART X. EXPONENTIAL
EQUATIONS}\label{part-x.-exponential-equations}

\section{38. Same-base method}\label{same-base-method}

If:

\[
a^{f(x)}=a^{g(x)}
\]

with a valid positive base \(a\neq1\), then:

\[
f(x)=g(x).
\]

Example:

\[
3^{x^2-4}=27^{x+2}.
\]

Rewrite:

\[
27=3^3.
\]

Thus:

\[
3^{x^2-4}
=
3^{3(x+2)}.
\]

Equate exponents:

\[
x^2-4=3x+6.
\]

\[
x^2-3x-10=0.
\]

Factor:

\[
(x-5)(x+2)=0.
\]

\[
\boxed{x=5,-2}.
\]

\begin{center}\rule{0.5\linewidth}{0.5pt}\end{center}

\section{\texorpdfstring{39. Exponential
\(u\)-substitution}{39. Exponential u-substitution}}\label{exponential-u-substitution}

Example:

\[
4^x-2^{x+1}=8.
\]

Rewrite:

\[
4^x=(2^x)^2.
\]

and:

\[
2^{x+1}=2(2^x).
\]

Let:

\[
u=2^x.
\]

Then:

\[
u^2-2u=8.
\]

\[
u^2-2u-8=0.
\]

\[
(u-4)(u+2)=0.
\]

So:

\[
u=4
\]

or

\[
u=-2.
\]

But:

\[
2^x>0
\]

for every real \(x\).

Reject \(-2\).

\[
2^x=4=2^2
\]

so:

\[
\boxed{x=2}.
\]

\begin{center}\rule{0.5\linewidth}{0.5pt}\end{center}

\section{PART XI. NONLINEAR SYSTEMS}\label{part-xi.-nonlinear-systems}

The standard strategy in your material is to substitute the linear
equation into the nonlinear equation, then solve and back-substitute.

\section{40. Line + parabola}\label{line-parabola}

Example:

\[
y=2x+1
\]

and

\[
y=x^2-2x+3.
\]

Since both equal \(y\):

\[
2x+1=x^2-2x+3.
\]

Move everything:

\[
x^2-4x+2=0.
\]

Solve:

\[
x=2\pm\sqrt2.
\]

Then plug each \(x\) into

\[
y=2x+1.
\]

That gives the two intersection points.

\begin{center}\rule{0.5\linewidth}{0.5pt}\end{center}

\section{41. Counting intersections using
discriminant}\label{counting-intersections-using-discriminant}

Suppose:

\[
y=mx+d
\]

and:

\[
y=ax^2+bx+c.
\]

Set equal:

\[
mx+d=ax^2+bx+c.
\]

Rearrange:

\[
ax^2+(b-m)x+(c-d)=0.
\]

Then use the discriminant.

\[
D>0\Rightarrow2\text{ intersections}
\]

\[
D=0\Rightarrow1\text{ intersection, tangent}
\]

\[
D<0\Rightarrow0\text{ intersections}.
\]

This is extremely SAT-friendly.

\begin{center}\rule{0.5\linewidth}{0.5pt}\end{center}

\section{42. Tangency parameter
problem}\label{tangency-parameter-problem}

Find \(k\) so that:

\[
y=2x+k
\]

and

\[
y=x^2-3x+1
\]

intersect exactly once.

Set equal:

\[
2x+k=x^2-3x+1.
\]

\[
x^2-5x+(1-k)=0.
\]

Exactly once means:

\[
D=0.
\]

\[
(-5)^2-4(1)(1-k)=0.
\]

\[
25-4+4k=0.
\]

\[
21+4k=0.
\]

\[
\boxed{k=-\frac{21}{4}}.
\]

\begin{center}\rule{0.5\linewidth}{0.5pt}\end{center}

\section{43. Quadratic + quadratic}\label{quadratic-quadratic}

Example:

\[
y=x^2-4
\]

\[
y=-x^2+4.
\]

Set equal:

\[
x^2-4=-x^2+4.
\]

\[
2x^2=8.
\]

\[
x^2=4.
\]

\[
x=\pm2.
\]

For either value:

\[
y=0.
\]

Solutions:

\[
\boxed{(2,0),(-2,0)}.
\]

\begin{center}\rule{0.5\linewidth}{0.5pt}\end{center}

\section{44. Absolute-value systems}\label{absolute-value-systems}

Example:

\[
y=2x-1
\]

\[
y=|x+2|.
\]

You can graph this instantly on Desmos.

Algebraically, split based on:

\[
x+2\ge0
\]

or

\[
x+2<0.
\]

The key is that after solving a branch, the answer must satisfy that
branch's condition.

\begin{center}\rule{0.5\linewidth}{0.5pt}\end{center}

\section{PART XII. REARRANGING NONLINEAR
FORMULAS}\label{part-xii.-rearranging-nonlinear-formulas}

Your document treats this as another tested skill: identify the target
variable and undo operations in reverse order.

\section{45. Simple square formula}\label{simple-square-formula}

\[
A=\pi r^2.
\]

Solve for \(r\):

\[
r^2=\frac A\pi.
\]

\[
r=\pm\sqrt{\frac A\pi}.
\]

But if \(r\) is a geometric radius:

\[
\boxed{r=\sqrt{\frac A\pi}}.
\]

Context removes the negative answer.

\begin{center}\rule{0.5\linewidth}{0.5pt}\end{center}

\section{46. Radical formula}\label{radical-formula}

\[
T=2\pi\sqrt{\frac Lg}.
\]

Divide:

\[
\frac{T}{2\pi}
=
\sqrt{\frac Lg}.
\]

Square:

\[
\frac{T^2}{4\pi^2}
=
\frac Lg.
\]

Multiply by \(g\):

\[
\boxed{L=\frac{gT^2}{4\pi^2}}.
\]

\begin{center}\rule{0.5\linewidth}{0.5pt}\end{center}

\section{47. Formula with fractions}\label{formula-with-fractions}

\[
\frac1R=\frac1{R_1}+\frac1{R_2}.
\]

Solve for \(R_1\).

Multiply everything by:

\[
RR_1R_2.
\]

Then:

\[
R_1R_2=RR_2+RR_1.
\]

Collect \(R_1\):

\[
R_1R_2-RR_1=RR_2.
\]

Factor:

\[
R_1(R_2-R)=RR_2.
\]

Thus:

\[
\boxed{R_1=\frac{RR_2}{R_2-R}}.
\]

\begin{center}\rule{0.5\linewidth}{0.5pt}\end{center}

\section{PART XIII. FUNCTIONS: DOMAIN AND
RANGE}\label{part-xiii.-functions-domain-and-range}

\section{48. Domain}\label{domain}

Domain means:

\[
\boxed{\text{allowed }x\text{-values}}.
\]

The SAT restrictions you care about most are:

{\def\LTcaptype{none} % do not increment counter
\begin{longtable}[]{@{}ll@{}}
\toprule\noalign{}
Expression & Restriction \\
\midrule\noalign{}
\endhead
\bottomrule\noalign{}
\endlastfoot
\(\sqrt{A}\) & \(A\ge0\) \\
\(\frac1B\) & \(B\neq0\) \\
real-life time/count/etc. & obey context \\
\end{longtable}
}

Example:

\[
f(x)=\frac{\sqrt{x-5}}{x-8}.
\]

Square root requires:

\[
x-5\ge0
\]

\[
x\ge5.
\]

Denominator requires:

\[
x\neq8.
\]

Domain:

\[
\boxed{[5,8)\cup(8,\infty)}.
\]

\begin{center}\rule{0.5\linewidth}{0.5pt}\end{center}

\section{49. Range}\label{range}

Range means:

\[
\boxed{\text{possible }y\text{-values}}.
\]

Example:

\[
f(x)=3+\sqrt{2x-4}.
\]

A square root is always nonnegative:

\[
\sqrt{2x-4}\ge0.
\]

Therefore:

\[
f(x)\ge3.
\]

Range:

\[
\boxed{[3,\infty)}.
\]

\begin{center}\rule{0.5\linewidth}{0.5pt}\end{center}

\section{PART XIV. MAXIMUM AND
MINIMUM}\label{part-xiv.-maximum-and-minimum}

\section{50. Vertex formula}\label{vertex-formula}

For:

\[
f(x)=ax^2+bx+c,
\]

the vertex \(x\)-coordinate is:

\[
\boxed{x_v=-\frac{b}{2a}}.
\]

Then:

\[
\boxed{y_v=f(x_v)}.
\]

If \(a>0\):

minimum.

If \(a<0\):

maximum.

\begin{center}\rule{0.5\linewidth}{0.5pt}\end{center}

\section{51. Example}\label{example}

\[
f(x)=-2x^2+12x-7.
\]

\[
a=-2,\quad b=12.
\]

\[
x_v=-\frac{12}{2(-2)}
=
3.
\]

Now:

\[
f(3)
=
-2(9)+36-7
=
11.
\]

Because the parabola opens downward:

\[
\boxed{\text{maximum}=11}
\]

at:

\[
\boxed{x=3}.
\]

Be careful about the wording.

``What is the maximum value?'' → \(11\)

``At what value of \(x\) does the maximum occur?'' → \(3\)

\begin{center}\rule{0.5\linewidth}{0.5pt}\end{center}

\section{PART XV. INCREASING AND
DECREASING}\label{part-xv.-increasing-and-decreasing}

For an upward parabola with vertex \(x=h\):

\[
\boxed{\text{decreasing on }(-\infty,h)}
\]

\[
\boxed{\text{increasing on }(h,\infty)}.
\]

For a downward parabola, reverse them.

Example:

\[
f(x)=x^2-4x+1.
\]

Vertex:

\[
x=-\frac{-4}{2}=2.
\]

Opens upward.

So:

\[
\boxed{\text{decreasing on }(-\infty,2)}
\]

and

\[
\boxed{\text{increasing on }(2,\infty)}.
\]

\begin{center}\rule{0.5\linewidth}{0.5pt}\end{center}

\section{PART XVI. POLYNOMIAL END
BEHAVIOR}\label{part-xvi.-polynomial-end-behavior}

Only the \textbf{leading term} matters.

Suppose:

\[
f(x)=ax^n+\cdots
\]

There are four possibilities.

{\def\LTcaptype{none} % do not increment counter
\begin{longtable}[]{@{}llll@{}}
\toprule\noalign{}
Degree & Leading coefficient & Left end & Right end \\
\midrule\noalign{}
\endhead
\bottomrule\noalign{}
\endlastfoot
Even & positive & up & up \\
Even & negative & down & down \\
Odd & positive & down & up \\
Odd & negative & up & down \\
\end{longtable}
}

Think:

\subsubsection{Even degree}\label{even-degree}

ends go the \textbf{same direction}.

\subsubsection{Odd degree}\label{odd-degree}

ends go \textbf{opposite directions}.

Then the sign of \(a\) tells you the right-hand end.

Example:

\[
f(x)=-3x^4+100x^3-5.
\]

Highest degree:

\[
-3x^4.
\]

Even + negative:

\[
\boxed{\text{both ends down}}.
\]

The \(100x^3\) does not control long-run behavior.

\begin{center}\rule{0.5\linewidth}{0.5pt}\end{center}

\section{PART XVII. GRAPH FEATURES OF RATIONAL
FUNCTIONS}\label{part-xvii.-graph-features-of-rational-functions}

\section{52. Intercepts}\label{intercepts}

\(x\)-intercept:

\[
f(x)=0.
\]

For a rational function, this usually means numerator \(=0\), provided
denominator \(\neq0\).

\(y\)-intercept:

\[
f(0).
\]

\begin{center}\rule{0.5\linewidth}{0.5pt}\end{center}

\section{53. Vertical asymptotes}\label{vertical-asymptotes}

For:

\[
f(x)=\frac{N(x)}{D(x)},
\]

a vertical asymptote occurs when:

\[
D(x)=0
\]

after simplification, provided the corresponding factor did not cancel.

Example:

\[
\frac{2x+1}{(x-3)(x+2)}.
\]

VA:

\[
\boxed{x=3,\;x=-2}.
\]

\begin{center}\rule{0.5\linewidth}{0.5pt}\end{center}

\section{54. Holes vs vertical
asymptotes}\label{holes-vs-vertical-asymptotes}

Example:

\[
f(x)=
\frac{(x-2)(x+3)}
{(x-2)(x+5)}.
\]

Cancel:

\[
f(x)=\frac{x+3}{x+5},
\]

but original restrictions are still:

\[
x\neq2,-5.
\]

At \(x=2\), the factor canceled.

So:

\[
\boxed{x=2\text{ is a hole}}.
\]

At \(x=-5\), denominator remains zero.

So:

\[
\boxed{x=-5\text{ is a vertical asymptote}}.
\]

Memorize:

\[
\boxed{\text{cancelled denominator factor} \rightarrow \text{hole}}
\]

\[
\boxed{\text{uncancelled denominator zero} \rightarrow \text{VA}}
\]

\begin{center}\rule{0.5\linewidth}{0.5pt}\end{center}

\section{55. Horizontal asymptotes}\label{horizontal-asymptotes}

For:

\[
\frac{\text{polynomial}}{\text{polynomial}},
\]

compare degrees.

\subsubsection{Numerator degree \textless{} denominator
degree}\label{numerator-degree-denominator-degree}

\[
\boxed{y=0}.
\]

Example:

\[
\frac{x}{x^2+1}.
\]

\subsubsection{Degrees equal}\label{degrees-equal}

Ratio of leading coefficients.

\[
\boxed{
y=
\frac{\text{leading numerator coefficient}}
{\text{leading denominator coefficient}}
}
\]

Example:

\[
\frac{5x^3-2x}{2x^3+7}.
\]

\[
\boxed{y=\frac52}.
\]

\subsubsection{Numerator degree \textgreater{} denominator
degree}\label{numerator-degree-denominator-degree-1}

No ordinary horizontal asymptote.

\begin{center}\rule{0.5\linewidth}{0.5pt}\end{center}

\section{PART XVIII. FUNCTION
TRANSFORMATIONS}\label{part-xviii.-function-transformations}

You need these cold.

{\def\LTcaptype{none} % do not increment counter
\begin{longtable}[]{@{}ll@{}}
\toprule\noalign{}
Expression & Transformation \\
\midrule\noalign{}
\endhead
\bottomrule\noalign{}
\endlastfoot
\(f(x)+k\) & up \(k\) \\
\(f(x)-k\) & down \(k\) \\
\(f(x-h)\) & right \(h\) \\
\(f(x+h)\) & left \(h\) \\
\(-f(x)\) & reflect over \(x\)-axis \\
\(f(-x)\) & reflect over \(y\)-axis \\
\(af(x), a>1\) & vertical stretch \\
\(af(x),0<a<1\) & vertical compression \\
\end{longtable}
}

The horizontal shift is the famous trap.

\[
f(x-5)
\]

moves \textbf{right 5}.

\[
f(x+5)
\]

moves \textbf{left 5}.

Inside is opposite.

Outside is intuitive.

\begin{center}\rule{0.5\linewidth}{0.5pt}\end{center}

\section{56. Transformation example}\label{transformation-example}

Parent:

\[
f(x)=\sqrt x.
\]

Shift right \(2\):

\[
\sqrt{x-2}.
\]

Reflect over \(x\)-axis:

\[
-\sqrt{x-2}.
\]

Shift down \(5\):

\[
\boxed{-\sqrt{x-2}-5}.
\]

\begin{center}\rule{0.5\linewidth}{0.5pt}\end{center}

\section{PART XIX. POLYNOMIAL ROOT
MULTIPLICITY}\label{part-xix.-polynomial-root-multiplicity}

\section{57. Multiplicity tells you graph
behavior}\label{multiplicity-tells-you-graph-behavior}

Suppose a polynomial contains:

\[
(x-r)^m.
\]

Then \(r\) is a root with multiplicity \(m\).

If multiplicity is odd:

\[
\boxed{\text{graph crosses}}
\]

If multiplicity is even:

\[
\boxed{\text{graph touches/bounces}}.
\]

Examples:

\[
(x-2)
\]

multiplicity \(1\): crosses.

\[
(x-2)^2
\]

multiplicity \(2\): touches.

\[
(x-2)^3
\]

multiplicity \(3\): crosses, usually flatter/S-shaped.

\begin{center}\rule{0.5\linewidth}{0.5pt}\end{center}

\section{58. Full polynomial graph
example}\label{full-polynomial-graph-example}

\[
f(x)=-(x+1)^2(x-3)(x-5)^3.
\]

Roots:

\[
x=-1,3,5.
\]

At \(-1\):

multiplicity \(2\) → bounce.

At \(3\):

multiplicity \(1\) → cross.

At \(5\):

multiplicity \(3\) → cross.

Degree:

\[
2+1+3=6.
\]

Notice here the leading sign is negative, so even degree + negative
coefficient means both ends down.

When doing these questions, calculate degree from \textbf{multiplicity},
not just number of distinct roots.

\begin{center}\rule{0.5\linewidth}{0.5pt}\end{center}

\section{59. Least possible degree}\label{least-possible-degree}

Graph:

cross at \(x=-1\)

touch at \(x=2\)

cross at \(x=5\).

Minimum multiplicities:

\[
1+2+1=4.
\]

So minimum degree:

\[
\boxed4.
\]

A touch cannot come from multiplicity \(1\).

\begin{center}\rule{0.5\linewidth}{0.5pt}\end{center}

\section{PART XX. EXPONENTIAL
FUNCTIONS}\label{part-xx.-exponential-functions}

This is extremely important on the Digital SAT.

\section{60. General form}\label{general-form}

\[
\boxed{f(t)=a(b)^t}.
\]

\(a\):

\[
\boxed{\text{initial value}}
\]

because:

\[
f(0)=ab^0=a.
\]

\(b\):

\[
\boxed{\text{multiplicative factor per time period}}.
\]

\begin{center}\rule{0.5\linewidth}{0.5pt}\end{center}

\section{61. Growth}\label{growth}

Increase by \(r\%\):

\[
\boxed{b=1+\frac r{100}}.
\]

8\% increase:

\[
b=1.08.
\]

15\% increase:

\[
b=1.15.
\]

So:

\[
500(1.08)^t
\]

means initial value \(500\), increasing \(8\%\) each period.

\begin{center}\rule{0.5\linewidth}{0.5pt}\end{center}

\section{62. Decay}\label{decay}

Decrease by \(r\%\):

\[
\boxed{b=1-\frac r{100}}.
\]

12\% depreciation:

\[
1-0.12=0.88.
\]

If a \$25,000 car loses 12\% annually:

\[
\boxed{V(t)=25000(0.88)^t}.
\]

Do not use \(0.12\).

The car retains \(88\%\), not \(12\%\).

\begin{center}\rule{0.5\linewidth}{0.5pt}\end{center}

\section{63. Multiplier ↔ percentage
conversion}\label{multiplier-percentage-conversion}

If multiplier:

\[
1.27,
\]

percentage growth:

\[
(1.27-1)(100)=27\%.
\]

If multiplier:

\[
0.84,
\]

percentage decrease:

\[
(1-0.84)(100)=16\%.
\]

\begin{center}\rule{0.5\linewidth}{0.5pt}\end{center}

\section{64. Doubling time}\label{doubling-time}

If a quantity doubles every \(d\) units:

\[
\boxed{f(t)=a2^{t/d}}.
\]

Example: initial population \(300\), doubles every 5 hours:

\[
\boxed{P(t)=300(2)^{t/5}}.
\]

At \(t=5\):

\[
300(2)^1=600.
\]

At \(t=10\):

\[
300(2)^2=1200.
\]

\begin{center}\rule{0.5\linewidth}{0.5pt}\end{center}

\section{65. Half-life}\label{half-life}

If something halves every \(h\) units:

\[
\boxed{f(t)=a\left(\frac12\right)^{t/h}}.
\]

50 grams, half-life 30 days:

\[
\boxed{M(t)=50\left(\frac12\right)^{t/30}}.
\]

At 30 days:

\[
25.
\]

At 60:

\[
12.5.
\]

\begin{center}\rule{0.5\linewidth}{0.5pt}\end{center}

\section{66. Scaled time, one of the nastiest SAT
traps}\label{scaled-time-one-of-the-nastiest-sat-traps}

Suppose:

\[
A(t)=18000(1.1025)^{t/2}.
\]

Do \textbf{not} say it grows 10.25\% every year.

The exponent is:

\[
\frac t2.
\]

So one full multiplication by \(1.1025\) occurs every \textbf{2 years}.

Thus:

\[
\boxed{10.25\%\text{ growth every 2 years}}.
\]

What is the per-year multiplier?

\[
(1.1025)^{1/2}.
\]

Since:

\[
1.05^2=1.1025,
\]

the annual factor is:

\[
1.05.
\]

So:

\[
\boxed{5\%\text{ growth per year}}.
\]

\begin{center}\rule{0.5\linewidth}{0.5pt}\end{center}

\section{67. Growth over several
periods}\label{growth-over-several-periods}

Suppose something grows \(8\%\) every 2 years.

What percent does it grow over 6 years?

Number of periods:

\[
\frac62=3.
\]

Overall multiplier:

\[
1.08^3.
\]

\[
1.08^3=1.259712.
\]

Growth percentage:

\[
(1.259712-1)(100)
=
25.9712\%.
\]

Approximately:

\[
\boxed{26.0\%}.
\]

You do not simply calculate:

\[
8\%\times3=24\%
\]

because exponential growth compounds.

\begin{center}\rule{0.5\linewidth}{0.5pt}\end{center}

\section{68. Exponential asymptotes}\label{exponential-asymptotes}

For:

\[
f(x)=ab^x+k,
\]

horizontal asymptote:

\[
\boxed{y=k}.
\]

Example:

\[
f(x)=5(2)^x-7.
\]

HA:

\[
\boxed{y=-7}.
\]

For \(a>0\), range:

\[
\boxed{(-7,\infty)}.
\]

\begin{center}\rule{0.5\linewidth}{0.5pt}\end{center}

\section{PART XXI. RATIONAL FUNCTION
TRANSFORMATIONS}\label{part-xxi.-rational-function-transformations}

\section{69. Parent function}\label{parent-function}

\[
\boxed{f(x)=\frac1x}.
\]

Transformed form:

\[
\boxed{f(x)=\frac a{x-h}+k}.
\]

Immediately:

\[
\boxed{\text{VA}:x=h}
\]

\[
\boxed{\text{HA}:y=k}.
\]

Example:

\[
f(x)=\frac3{x-2}+1.
\]

VA:

\[
x=2.
\]

HA:

\[
y=1.
\]

Domain:

\[
x\neq2.
\]

Range:

\[
y\neq1.
\]

\begin{center}\rule{0.5\linewidth}{0.5pt}\end{center}

\section{PART XXII. DESMOS: WHAT YOU SHOULD ACTUALLY USE IT
FOR}\label{part-xxii.-desmos-what-you-should-actually-use-it-for}

Because the Digital SAT includes Desmos, you should know when it saves
time.

\subsection{Roots}\label{roots}

Graph:

\[
y=f(x).
\]

The \(x\)-intercepts are the real roots.

\subsection{Systems}\label{systems}

Graph both equations.

Each intersection is one solution.

This is particularly useful for ``how many solutions?'' questions.

\subsection{Vertex}\label{vertex}

Graph the quadratic and select the turning point.

Useful for min/max questions when exact symbolic manipulation is
unnecessary.

\subsection{Equivalent expressions}\label{equivalent-expressions}

Graph:

\[
y=\text{original}
\]

and:

\[
y=\text{answer choice}.
\]

If they overlap for the relevant domain, they are equivalent.

Be careful about \textbf{holes/domain restrictions}.

\subsection{Regression from a table}\label{regression-from-a-table}

Quadratic:

\[
y_1\sim ax_1^2+bx_1+c.
\]

Exponential:

\[
y_1\sim ab^{x_1}.
\]

These can recover models rapidly.

\subsection{When you should NOT automatically use
Desmos}\label{when-you-should-not-automatically-use-desmos}

Simple factoring such as:

\[
x^2-9.
\]

Simple coefficient interpretation such as:

\[
300(1.07)^t.
\]

Exact radical answers such as:

\[
\frac{7+\sqrt{73}}4.
\]

Basic exponent rules.

You need enough algebra fluency that opening Desmos does not make a
10-second question take a minute.

\begin{center}\rule{0.5\linewidth}{0.5pt}\end{center}

\section{PART XXIII. THE TRAPS YOU NEED TO BE PARANOID
ABOUT}\label{part-xxiii.-the-traps-you-need-to-be-paranoid-about}

Your source explicitly emphasizes these traps.

{\def\LTcaptype{none} % do not increment counter
\begin{longtable}[]{@{}ll@{}}
\toprule\noalign{}
Trap & Correct idea \\
\midrule\noalign{}
\endhead
\bottomrule\noalign{}
\endlastfoot
\((a+b)^2=a^2+b^2\) & Wrong. Include \(2ab\) \\
\(x^2+49\) factors like diff. squares & No real factorization \\
Radical equation candidates & Check original \\
Cancelled denominator & Original excluded value remains excluded \\
\(|x|=-5\) & No solution \\
\(f(x+3)\) & Shift left 3 \\
\(D=0\), then \(k^2=36\) & \(k=\pm6\) \\
Even multiplicity & Bounce \\
Odd multiplicity & Cross \\
\(1.08\) & 8\% growth \\
\(0.92\) & 8\% decay \\
Completing square with \(a\neq1\) & Factor \(a\) first \\
Rational simplification & Cancel factors, not terms \\
\(b^x\) for positive base & Always positive \\
\(1.08^{t/2}\) & 8\% every 2 units \\
\end{longtable}
}

\begin{center}\rule{0.5\linewidth}{0.5pt}\end{center}

\section{PART XXIV. THE MASTER SAT METHOD
CHECKLIST}\label{part-xxiv.-the-master-sat-method-checklist}

This is the checklist I would actually memorize.

{\def\LTcaptype{none} % do not increment counter
\begin{longtable}[]{@{}
  >{\raggedright\arraybackslash}p{(\linewidth - 2\tabcolsep) * \real{0.5000}}
  >{\raggedright\arraybackslash}p{(\linewidth - 2\tabcolsep) * \real{0.5000}}@{}}
\toprule\noalign{}
\begin{minipage}[b]{\linewidth}\raggedright
If you see\ldots{}
\end{minipage} & \begin{minipage}[b]{\linewidth}\raggedright
Do this
\end{minipage} \\
\midrule\noalign{}
\endhead
\bottomrule\noalign{}
\endlastfoot
``Equivalent expression'' & factor / expand / exponent rules / test
values \\
Polynomial with 2 terms & GCF → diff squares → cubes \\
Polynomial with 3 terms & GCF → perfect square → trinomial \\
Polynomial with 4 terms & GCF → grouping \\
Rational expression & factor numerator + denominator first \\
Square root + exponents & convert to fractional exponents \\
Quadratic roots & factoring / square root / quadratic formula \\
``How many real solutions?'' quadratic & discriminant \\
Radical equation & isolate → square → solve → check \\
Rational equation & restrictions → LCD → solve → reject forbidden
values \\
Absolute-value equation & isolate → positive and negative cases \\
Exponential equation & same base or \(u=b^x\) \\
Line + parabola system & substitute → quadratic \\
``Exactly one intersection'' & discriminant \(=0\) \\
Quadratic min/max & vertex \\
Quadratic roots & factored form \\
Quadratic \(y\)-intercept & standard form \\
Domain & square roots \(\ge0\); denominators \(\neq0\) \\
Range & reason from graph/output restrictions \\
Polynomial end behavior & leading term only \\
Rational VA & uncancelled denominator zero \\
Rational hole & cancelled denominator factor \\
Exponential model & initial ×
multiplier\(^{\text{number of periods}}\) \\
Doubling & \(a2^{t/d}\) \\
Half-life & \(a(1/2)^{t/h}\) \\
\(f(x-h)\) & right \\
\(f(x+h)\) & left \\
even multiplicity & touch \\
odd multiplicity & cross \\
\end{longtable}
}

\begin{center}\rule{0.5\linewidth}{0.5pt}\end{center}

\section{PART XXV. FORMULA SHEET TO
MEMORIZE}\label{part-xxv.-formula-sheet-to-memorize}

\subsection{Factoring}\label{factoring}

\[
\boxed{a^2-b^2=(a-b)(a+b)}
\]

\[
\boxed{a^2+2ab+b^2=(a+b)^2}
\]

\[
\boxed{a^2-2ab+b^2=(a-b)^2}
\]

\[
\boxed{a^3+b^3=(a+b)(a^2-ab+b^2)}
\]

\[
\boxed{a^3-b^3=(a-b)(a^2+ab+b^2)}
\]

\subsection{Quadratics}\label{quadratics}

\[
\boxed{x=\frac{-b\pm\sqrt{b^2-4ac}}{2a}}
\]

\[
\boxed{D=b^2-4ac}
\]

\[
\boxed{x_v=-\frac b{2a}}
\]

\[
\boxed{y_v=f(x_v)}
\]

\subsection{Vertex form}\label{vertex-form-1}

\[
\boxed{f(x)=a(x-h)^2+k}
\]

vertex:

\[
\boxed{(h,k)}.
\]

\subsection{Factored form}\label{factored-form-1}

\[
\boxed{f(x)=a(x-r_1)(x-r_2)}
\]

roots:

\[
\boxed{r_1,r_2}.
\]

axis:

\[
\boxed{x=\frac{r_1+r_2}{2}}.
\]

\subsection{Completing square}\label{completing-square}

\[
\boxed{x^2+bx=
\left(x+\frac b2\right)^2-
\left(\frac b2\right)^2}
\]

\subsection{Fractional exponents}\label{fractional-exponents}

\[
\boxed{a^{m/n}=\sqrt[n]{a^m}}
\]

\subsection{Exponential growth}\label{exponential-growth}

\[
\boxed{A=A_0(1+r)^t}
\]

\subsection{Exponential decay}\label{exponential-decay}

\[
\boxed{A=A_0(1-r)^t}
\]

\subsection{\texorpdfstring{Change every \(d\)
units}{Change every d units}}\label{change-every-d-units}

\[
\boxed{A=A_0b^{t/d}}
\]

\subsection{Doubling}\label{doubling}

\[
\boxed{A=A_0\,2^{t/d}}
\]

\subsection{Half-life}\label{half-life-1}

\[
\boxed{A=A_0\left(\frac12\right)^{t/h}}
\]

\subsection{Compound interest}\label{compound-interest}

\[
\boxed{
A=P\left(1+\frac rn\right)^{nt}
}
\]

\subsection{Rational horizontal
asymptote}\label{rational-horizontal-asymptote}

\[
\boxed{
\deg N<\deg D\Rightarrow y=0
}
\]

\[
\boxed{
\deg N=\deg D
\Rightarrow
y=\frac{\text{leading coeff. }N}
{\text{leading coeff. }D}
}
\]

\subsection{Rational transformation}\label{rational-transformation}

\[
\boxed{f(x)=\frac a{x-h}+k}
\]

\[
\boxed{\text{VA }x=h,\quad\text{HA }y=k}
\]

\subsection{Absolute value}\label{absolute-value}

\[
\boxed{|A|=B\Rightarrow A=B\text{ or }A=-B}
\]

provided \(B\ge0\).

\begin{center}\rule{0.5\linewidth}{0.5pt}\end{center}

\section{Final ``am I SAT-ready?''
checklist}\label{final-am-i-sat-ready-checklist}

You should be able to do every one of these without needing notes:

\begin{enumerate}
\def\labelenumi{\arabic{enumi}.}
\tightlist
\item
  Expand \((3x-5)^2\) without forgetting the middle term.
\item
  Factor simple and hard trinomials.
\item
  Recognize difference of squares and perfect squares immediately.
\item
  Simplify rational expressions while preserving domain restrictions.
\item
  Convert radicals into rational exponents.
\item
  Complete the square, including when \(a\neq1\).
\item
  Switch mentally among standard, factored, and vertex quadratic forms.
\item
  Choose the fastest quadratic-solving method.
\item
  Use \(D=b^2-4ac\) for 0/1/2 real roots and tangency.
\item
  Recognize \(x^4+bx^2+c\) as a quadratic in \(x^2\).
\item
  Solve radical equations and check for extraneous roots.
\item
  Solve rational equations and reject denominator-zero candidates.
\item
  Split absolute-value equations correctly.
\item
  Rewrite exponential equations to a common base.
\item
  Use \(u=b^x\) for quadratic-in-exponential equations.
\item
  Solve/count intersections in nonlinear systems.
\item
  Rearrange formulas containing squares, roots, and fractions.
\item
  Determine domain and range.
\item
  Find quadratic min/max and know whether the question wants \(x\) or
  \(y\).
\item
  Determine increasing/decreasing intervals from the vertex.
\item
  Determine polynomial end behavior from degree parity and leading sign.
\item
  Identify holes versus vertical asymptotes.
\item
  Identify horizontal asymptotes.
\item
  Interpret transformations, especially horizontal shifts.
\item
  Read root multiplicity from factors and graph behavior.
\item
  Translate percent growth/decay into multipliers.
\item
  Handle half-life and doubling-time models.
\item
  Interpret scaled exponents such as \(t/3\).
\item
  Compute compound percentage change rather than adding percentages.
\item
  Know when Desmos is faster and when paper algebra is faster.
\end{enumerate}

If you can do those 30 things confidently, you have covered the
substantive material in the B2 Advanced Math document rather than merely
memorizing its formulas. The file itself frames the same domain around
nonlinear relationships involving quadratics, exponentials, polynomials,
radicals, rationals, and systems.

\end{document}
